\documentclass{includes/Thesis}

\addbibresource{parts/refs.bib}
\usepackage{float}
\usepackage{placeins}
\usetikzlibrary{arrows.meta,backgrounds,fit,positioning}

\definecolor{listingbg}{gray}{0.96}
\definecolor{listingkw}{RGB}{0,74,140}
\definecolor{listingtype}{RGB}{120,78,0}
\definecolor{listingconst}{RGB}{90,40,120}
\definecolor{listingcomment}{RGB}{0,120,60}
\definecolor{listingstring}{RGB}{163,21,21}
\renewcommand{\lstlistingname}{Листинг}
\DeclareCaptionLabelSeparator{hyphen}{ - }
\captionsetup{labelsep=hyphen}
\emergencystretch=2em

\let\oldtexttt\texttt
\renewcommand{\texttt}[1]{\mbox{\oldtexttt{#1}}}
\newcounter{appendixlstbase}

\lstdefinelanguage{PseudoPy}{
    language=Python,
    morekeywords={
        func,emit,return,copy,spawn_partition_generators,
        generate_sentence,resolve_selected_write_strategy,
        generate_partition_batches,build_text_pool,
        export_record_batch,register_views,deliver_batches,
        buffer_message,flush_ready_partitions,
        next_token,generate_verb_phrase,generate_noun_phrase,
        trim_trailing_space
    },
    morekeywords=[2]{
        TextPool,RowRandomInt,RowRandomLong,
        RandomBoundedInt,RandomBoundedLong,
        RandomText,TableBatch,ArrowArrayStream,
        RecordBatch,GenerationContext
    },
    morekeywords=[3]{
        INT32_MAX,LINE_COUNT_MAX,TEXT_POOL_SIZE
    },
    morecomment=[s]{/*}{*/},
    sensitive=true
}

\lstdefinelanguage{DistsDSS}{
    morekeywords={BEGIN,END,COUNT,begin,end,count},
    sensitive=false,
    morecomment=[l]{\#}
}

\lstdefinestyle{reportlisting}{
    style=reportlistinginline
}

\lstdefinestyle{reportlistinginline}{
    basicstyle=\ttfamily\footnotesize,
    breaklines=true,
    columns=fullflexible,
    showstringspaces=false,
    keepspaces=true,
    frame=single,
    framerule=0.4pt,
    rulecolor=\color{black},
    backgroundcolor=\color{listingbg},
    framesep=6pt,
    captionpos=b,
    aboveskip=6pt,
    belowskip=10pt,
    numbers=none
}

\lstdefinestyle{reportpseudocode}{
    style=reportlisting,
    language=PseudoPy,
    keywordstyle=\color{listingkw}\bfseries,
    keywordstyle=[2]\color{listingtype}\bfseries,
    keywordstyle=[3]\color{listingconst}\bfseries,
    commentstyle=\color{listingcomment}\itshape,
    stringstyle=\color{listingstring}
}

\lstdefinestyle{reportpseudocodeinline}{
    style=reportlistinginline,
    language=PseudoPy,
    keywordstyle=\color{listingkw}\bfseries,
    keywordstyle=[2]\color{listingtype}\bfseries,
    keywordstyle=[3]\color{listingconst}\bfseries,
    commentstyle=\color{listingcomment}\itshape,
    stringstyle=\color{listingstring}
}

\lstdefinestyle{reportdss}{
    style=reportlisting,
    language=DistsDSS,
    keywordstyle=\color{listingkw}\bfseries,
    commentstyle=\color{listingcomment}\itshape,
    stringstyle=\color{listingstring}
}

\begin{thesis}[
    title = Генератор данных для аналитического бенчмарка СУБД TPC-H,
    year = 2026,
    %
    authorGroup = БПИ225,
    authorName = М. А. Дергоусов,
    %
    academicTeacherTitle = {профессор департамента программной инженерии ФКН НИУ ВШЭ},
    academicTeacherName = В. В. Шилов,
    academicSupervisorTitle = {Академический руководитель\\образовательной программы\\<<Программная инженерия>>},
    academicSupervisorName = Н. А. Павлочев,
    %
    UDC = 004.65,
    %
    hasConsultant,
    consultantTitle = {разработчик Яндекс Технологии},
    consultantName = М. Т. Садуллаев,
    %
    keywordsRu = OLAP; СУБД; TPC-H; Apache Arrow; Parquet; ORC; DuckDB,
    keywordsEn = OLAP; DBMS; TPC-H; Apache Arrow; Parquet; ORC; DuckDB,
]
    \setAbstractResource{parts/abstract-ru}{parts/abstract-en}

    \setTerminologyResource{parts/terminology}

    \setIntroResource{parts/intro}

    \addChapter{Предметная область и анализ аналогов}{parts/chapter1}
    \addChapter{Архитектура генератора данных}{parts/chapter2}
    \addChapter{Реализация генератора данных}{parts/chapter3}
    \addAppendix{Генератор таблицы region}{parts/appendix-a}
    \addAppendix{Генератор таблицы nation}{parts/appendix-b}
    \addAppendix{Генератор таблицы supplier}{parts/appendix-c}
    \addAppendix{Генератор таблицы part}{parts/appendix-d}
    \addAppendix{Генератор таблицы partsupp}{parts/appendix-e}
    \addAppendix{Генератор таблицы customer}{parts/appendix-f}
    \addAppendix{Генератор таблицы orders}{parts/appendix-g}
    \addAppendix{Генератор таблицы lineitem}{parts/appendix-h}

    \setConclusionResource{parts/conclusion}
\end{thesis}
